%% appF-manifest \dash{} Manifest
%
% Three things, not one: the diagram manifest, the measurement table, and the
% outstanding-work ledgers -- printed rather than merely counted by the build,
% because a reader holding the paper cannot run `make debt`.
%
% EVERY NUMBER ON THESE PAGES IS COMPUTED. code/appf_ledgers.py counts them
% out of the tree and writes figures/values/appf.tex, so `make verify` fails
% when a ledger moves and this appendix does not. That is deliberate and it is
% the whole design: an appendix of hand-typed counts is precisely what went
% wrong elsewhere -- a debt ledger that said 27 where its own tool printed 54,
% a pages-per-frame ratio left at a superseded figure, seven wrong part ranges
% on page one -- all inside one week, and none of it visible to any gate.
%
% TWO THINGS ARE DELIBERATELY NOT NUMBERS HERE, and the reasons are in the
% script's docstring as well:
%
%   * how many of the ten experiments have run. A count of how many is a claim
%     about this book that nothing derives from anything, and notes/01 s17
%     records that "none has been run" sat on the page for five programs after
%     one had. The rows carry their own status; there is no total.
%
%   * the page-level ledgers (stranded openers, orphaned cues, orphan tails).
%     Those are properties of the TYPESETTING, and this book's two TeX
%     installations paginate differently, so a printed count would be a fact
%     about one machine wearing the clothes of a fact about the book.
%
% Counting convention: every figure in this appendix describes THE COPY YOU
% ARE HOLDING, which is one edition. `make debt` counts both and therefore
% prints twice several of these. Where the difference matters the appendix
% gives both numbers rather than leaving the reader to guess which they have.

\chapter{Manifest}
\label{app:F}

\noindent
Appendices~\ref{app:B}, \ref{app:C} and~\ref{app:D} answer the fair criticism
of a book made of frames: that it teaches well and is useless to look
something up in afterwards. This one answers a different question, and it is
the one a book is least inclined to ask itself.

\emph{What does this book still owe you?}

Every number below is counted out of the source by a script and pulled onto
the page mechanically, exactly as every other number in this book is. That is
not tidiness. An appendix of hand-typed counts goes stale in weeks and nothing
can see it happen \dash{} which is a claim this book can make with some
authority, having done it.

\medskip
\begin{note}
Every figure here describes the copy you are holding, which is one edition.
The build's own ledgers count both editions at once and so print twice several
of these; where that matters, both numbers are given.
\end{note}

\section{The diagrams}
\label{sec:F-diagrams}
\index{manifest!diagrams}

There are $\val{appf.diagrams}$ diagrams in this edition and
$\val{appf.diagrams.both}$ across the two, because each is drawn separately in
each language: a diagram with English labels in the Polish edition is the kind
of half-translation that makes a bilingual book feel machine-made.

The source of truth is a text file, committed, so a change to a diagram
reviews as a diff rather than as a new picture. The rendered image is build
output and is not kept.

\textbf{The list below is not maintained.} Each entry is written by the figure macro
itself at the point where the figure is placed, so a diagram cannot be in the
book and missing from this list, nor listed here and absent from the book. It
is the same mechanism as Appendix~\ref{app:C}, and it is used for the same
reason: a hand-kept index of a book's own contents is a claim about the book on
every line, and that is the one class of claim nothing here can check.

\listofdiagrams

\section{The measurements}
\label{sec:F-measurements}
\index{manifest!measurements}

This book's house rule is that a claim about behaviour \dash{} not about a
theorem \dash{} is demonstrated by a run rather than asserted. Ten experiments
were specified for that purpose. Each row below carries its own status.

\textbf{No total appears here, and the omission is the point.} A count of how
many have been run is a claim about this book that nothing derives from
anything, and it decays silently: the sentence \enquote{none has been run}
stayed on the page, and in the file the next author reads first, for five
programs after one had been.

\medskip
\textbf{Run, and reported in the program that ran them.}

\begin{center}
\begin{tabularx}{\linewidth}{@{}lXl@{}}
\toprule
\textbf{\#} & \textbf{What it measures} & \textbf{Where} \\
\midrule
E6 & Plain descent, momentum and Adam on a quadratic of known condition
  number; steps to tolerance against the predicted count &
  \ref{prog:P20} \\
E8 & Forward against reverse divergence fitted to one bimodal target;
  mode-covering against mode-seeking & \ref{prog:P30} \\
E9 & Score spread and attention entropy with and without the square-root
  divisor, across head sizes & \ref{prog:P25}, \ref{prog:P32} \\
E10 & A power-law fit on the shape the literature uses, with the fit's own
  extrapolation error reported & \ref{prog:P33} \\
\bottomrule
\end{tabularx}
\end{center}

\medskip
\textbf{Measured, but not claimed.} Each of these programs measured something
that closely resembles the specification beside it, and no pass ever claimed
the experiment. Whether the specification is met is a reading job on a written
program, and it is not an inference to draw from a table \dash{} which is why
the table says this rather than saying \enquote{run}.

\begin{center}
\begin{tabularx}{\linewidth}{@{}lXl@{}}
\toprule
\textbf{\#} & \textbf{What it specifies} & \textbf{Where} \\
\midrule
E1 & The score at which a naive $\softmax$ overflows each format, and what a
  non-maximal pivot costs & \ref{prog:P02} \\
E3 & The angle between random unit vectors as the dimension climbs, and the
  concentration towards a right angle & \ref{prog:P05} \\
E5 & Forward against reverse mode: time and peak memory against depth, and
  the measured cost of checkpointing & \ref{prog:P16} \\
\bottomrule
\end{tabularx}
\end{center}

\medskip
\textbf{Not run.}

\begin{center}
\begin{tabularx}{\linewidth}{@{}lXl@{}}
\toprule
\textbf{\#} & \textbf{What it would settle} & \textbf{Where} \\
\midrule
E2 & Hand-counted operations and bytes for one forward pass against a
  measured wall clock; where the model is wrong and by how much &
  \ref{prog:P03} \\
E4 & The singular-value spectrum of a real embedding matrix, and
  reconstruction error against rank \dash{} this one needs a trained model &
  \ref{prog:P11} \\
E7 & Interval width against evaluation-set size on a public benchmark; the
  size needed to resolve one point & \ref{prog:P27} \\
\bottomrule
\end{tabularx}
\end{center}

\section{What the build counts}
\label{sec:F-ledgers}
\index{manifest!ledgers}

Each of these is checked on every build, and each of them is a number this
book would rather you saw than took on trust.

\begin{center}
\begin{tabularx}{\linewidth}{@{}Xl@{}}
\toprule
\textbf{Ledger} & \textbf{Count} \\
\midrule
Programs written, of $\val{appf.programs}$ & $\val{appf.programs}$ \\
Teaching frames & $\val{appf.frames}$ \\
Frames the plan projected & $\val{appf.frames.planned}$ \\
Computed values, every one referenced and present & $\val{appf.values}$ \\
Console transcripts, every one produced by a script &
  $\val{appf.transcripts}$ \\
Diagrams & $\val{appf.diagrams}$ \\
Claims marked as unverified & $\val{appf.verifybox}$ \\
Frames that put a question to you &
  $\val{appf.elicit}$, or $\val{appf.elicit.pct}$ per cent \\
Polish renderings named in Appendix~\ref{app:D} & $\val{appf.terms.rend}$
  in $\val{appf.terms.rows}$ rows \\
Part ranges in the introduction, against the manifest & $\val{appf.parts}$ \\
\bottomrule
\end{tabularx}
\end{center}

\medskip
Two of those describe one edition where the build's own ledger describes both,
and the difference is a factor of two rather than a disagreement: there are
$\val{appf.transcripts.both}$ transcripts across the two editions and
$\val{appf.diagrams.both}$ diagrams, because each is produced separately in
each language. Quote whichever you mean and say which it is.

\medskip
Three more are checked on the finished page rather than in the source, and
they are deliberately not given a number here: whether a frame's rule and
margin number are ever stranded at the foot of a page with the frame overleaf,
whether a section heading is, and how often a frame's tail lands nearly alone
on a page. Those are properties of the typesetting, and the two typesetting
installations this book is built on break lines differently, so a count
printed here would be a fact about one machine dressed as a fact about the
book.

\section{What this book does not claim}
\label{sec:F-owed}
\index{manifest!outstanding}

\begin{warning}
\textbf{The 80/80 standard is not established, and this book may not claim
it.} The method it borrows was validated to that standard \dash{} eighty per
cent of readers reaching eighty per cent of the objectives \dash{} before its
own publication. That is a fact about the method and it is not transferable.
Nobody has been trialled on \emph{this} book. Until somebody has, the ledger
stays open, and the instrument when it happens is the scored Test exercises
rather than the Quiz: the Quiz runs on the Foundation programs only, and its
items serve on entry and on exit alike, so a score difference across it
measures memory of the questions.
\end{warning}

\medskip
Three other things are owed, and they are worth naming individually because
each is a different kind of debt.

\medskip
\textbf{The measurements that need a trained model.} Four claims in this book
are stated as mechanism and labelled as judgement rather than measured,
because settling them needs weights this book does not ship and does not
download: whether a real embedding matrix's spectrum falls the way low-rank
methods assume; whether the rank of a deep attention stack falls without any
narrow layer to force it; whether the basin a training walk lands in matters
at the scale people train at; and whether the independence the square-root
divisor assumes survives training. Each is marked where it appears.

\medskip
\textbf{The rate at which this book asks you questions is not even.} It runs
at roughly three quarters through the early Foundation programs and at about
half through the main sequence, and the book's own figure is
$\val{appf.elicit.pct}$ per cent. The gap has a cause rather than an excuse: a
program that derives a result has fewer places to stop and ask than one built
on numbers. It is still a real unevenness in the teaching rather than a
statistic, because a program that tells you more than it asks you is a program
you will retain less from, whatever it says about itself.

\medskip
\textbf{And the companion library is specified and unbuilt.} One stage per
part, every gradient checked against a finite difference, no accelerator
required. It exists as a design and as nothing else.

\medskip
That is the whole of what is outstanding, as far as anybody knows. The
qualification is not modesty: this appendix is itself a claim about the book,
of exactly the kind the preceding pages spend their length teaching you to
distrust, and the only reason to give it more weight than the others is that
its numbers are counted by a machine and its judgements are labelled as
judgements.

Which is all this book has ever asked you to require of anybody.
